\chapter{Reflexive Verben}
\label{cha:Reflexiv}

Reflexive Verben werden mit einem Reflexivpronomen gebraucht. Das Reflexivpronomen zeigt an, dass sich die Handlung auf das Subjekt des Satzes bezieht.

\section{Reflexivpronomen}
\label{sec:Reflexivpronomen}

\begin{center}
\begin{tabular}{|l|c|c|}
\hline
\textbf{Person} & \textbf{Akkusativ} & \textbf{Dativ} \\ \hline
ich & mich & mir \\ \hline
du & dich & dir \\ \hline
er/sie/es & sich & sich \\ \hline
wir & uns & uns \\ \hline
ihr & euch & euch \\ \hline
sie/Sie & sich & sich \\ \hline
\end{tabular}
\end{center}

\section{Arten von Reflexivverben}
\label{sec:ReflexivArten}

\subsection{Echte Reflexivverben}
Das Reflexivpronomen ist obligatorisch:
\begin{itemize}
    \item Ich wasche \textbf{mich}.
    \item Er setzt \textbf{sich}.
    \item Wir freuen \textbf{uns}.
    \item Sie entscheidet \textbf{sich}.
\end{itemize}

\subsection{Reflexiv mit Dativ}
Das reflexive Objekt steht im Dativ:
\begin{itemize}
    \item Ich wasche \textbf{mir} die Hände.
    \item Er kocht \textbf{sich} ein Ei.
    \item Sie nimmt \textbf{sich} Zeit.
\end{itemize}

\subsection{Reziproke Verben (wechselseitig)}
\begin{itemize}
    \item Wir sehen \textbf{einander}.
    \item Sie helfen \textbf{einander}.
    \item Sie verstehen \textbf{einander}.
\end{itemize}

\section{Häufige reflexive Verben}
\label{sec:ReflexivVerben}

\begin{itemize}
    \item \textbf{sich anmelden}: Ich melde mich an.
    \item \textbf{sich anziehen}: Er zieht sich an.
    \item \textbf{sich aufregen}: Sie regt sich auf.
    \item \textbf{sich beeilen}: Beeil dich!
    \item \textbf{sich beschweren}: Ich beschwere mich.
    \item \textbf{sich entscheiden}: Ich entscheide mich.
    \item \textbf{sich entschließen}: Er hat sich entschlossen.
    \item \textbf{sich erholen}: Wir erholen uns.
    \item \textbf{sich erkälten}: Ich habe mich erkältet.
    \item \textbf{sich erinnern}: Ich erinnere mich.
    \item \textbf{sich freuen}: Ich freue mich.
    \item \textbf{sich fühlen}: Ich fühle mich wohl.
    \item \textbf{sich interessieren}: Ich interessiere mich für Kunst.
    \item \textbf{sich konzentrieren}: Ich konzentriere mich.
    \item \textbf{sich langweilen}: Er langweilt sich.
    \item \textbf{sich leisten}: Das kann ich mir nicht leisten.
    \item \textbf{sich patschen}: Das Kind patscht sich.
    \item \textbf{sich streiten}: Sie streiten sich.
    \item \textbf{sich unterhalten}: Wir unterhalten uns.
    \item \textbf{sich verlieben}: Er hat sich verliebt.
    \item \textbf{sich verlaufen}: Ich habe mich verlaufen.
    \item \textbf{sich vorstellen}: Ich stelle mir etwas vor.
\end{itemize}

\section{Übungen zu reflexiven Verben}
\label{sec:ReflexivUebungen}

\begin{enumerate}
    \item Ich beschwere \underline{\hspace{1cm}} über die neue Waschmaschine. (mich)
    \item Die Gäste begrüßen \underline{\hspace{1cm}}. (sich)
    \item Ich sehe \underline{\hspace{1cm}} den Film nicht noch einmal an. (mir)
    \item Er interessiert \underline{\hspace{1cm}} für Kunst. (sich)
    \item Wir freuen \underline{\hspace{1cm}} auf den Urlaub. (uns)
    \item Du musst \underline{\hspace{1cm}} beeilen. (dich)
    \item Sie erinnert \underline{\hspace{1cm}} an den Termin. (sich)
    \item Ich kann \underline{\hspace{1cm}} das nicht leisten. (mir)
    \item Er bildet \underline{\hspace{1cm}} viel ein. (sich)
    \item Wir unterhalten \underline{\hspace{1cm}} gut. (uns)
\end{enumerate}

\section{Grammatik aktiv: Besonderheiten}
\label{sec:ReflexivBesonderheiten}

\begin{tcolorbox}[title=Wichtig]
Bei manchen Verben ändert sich die Bedeutung mit/reflexiv:
\end{tcolorbox}

\begin{itemize}
    \item \textbf{erinnern}: Ich erinnere ihn an den Termin. (= Ich sage ihm, dass...)
    \item \textbf{sich erinnern}: Ich erinnere mich an den Termin. (= Ich habe das Gedächtnis daran.)
    \item \textbf{treffen}: Ich treffe ihn. (= Wir begegnen uns.)
    \item \textbf{sich treffen}: Wir treffen uns. (= Wir haben eine Verabredung.)
\end{itemize}

\chapterend